% ARR/ACL-ready experiment and results section for the Final-197 matrix.
% Required in the main preamble: \usepackage{booktabs}
% Optional: rename section titles to match the surrounding manuscript.
% The RQ framing organizes the frozen six-cell experiment after completion;
% only the protocol, primary arm, and primary metrics were frozen in advance.

\section{Experiments}
\label{sec:experiments}

\subsection{Research Questions}
\label{sec:research-questions}

We organize the evaluation around four research questions.

\paragraph{RQ1: Cross-target effectiveness.}
How effectively does our target-grounded, multi-turn research dialogue elicit
goal-consistent responses from both a proprietary target (GPT-4o) and an
open-weight target (Qwen-2.5-7B-Instruct)?

\paragraph{RQ2: Dependence on verbatim goal exposure.}
Does the method require placing the original harmful request verbatim in the
target-visible persona, or can it operate when the target receives only
semantically related persona context and research dialogue?

\paragraph{RQ3: Contribution of pre-generated persona history.}
Does a goal-aware dialogue history improve elicitation beyond the dynamic
research dialogue itself?

\paragraph{RQ4: Robustness across safety evaluators.}
Do conclusions based on HarmBench remain directionally consistent under
Qwen3Guard-Gen-8B and Nemotron-3-Content-Safety?

The full-run prompt, seven-stage dialogue protocol, primary condition, and
primary metrics were frozen before evaluating all 197 cases.  The RQ framing
above organizes the resulting controlled matrix; we do not treat the
best-performing ablation as a retrospectively pre-registered primary method.

\subsection{Dataset}
\label{sec:final197-dataset}

We evaluate on \textsc{Final-197}, a conservatively filtered set of 197
explicit mental-health-related harmful requests.  The set contains 136 items
derived from the \emph{Between Help and Harm} collection, 22 from AILuminate
Suicide \& Self-Harm, 22 from BeaverTails, and 17 from Safety-Harmful.  The
original metadata marks 190 cases as H2 and seven as H1.  Selection required
an identifiable requested output and an explicit request for assistance with
the execution, preparation, optimization, concealment, or encouragement of
harm.  The filtering was LLM-assisted and did not involve an independent
clinical-expert consensus; we therefore treat this dataset as an evaluation
set for explicit harmful assistance rather than a clinical benchmark.

\subsection{Targets and Researcher Models}
\label{sec:models}

We use GPT-4o (\texttt{gpt-4o-2024-11-20}) and
Qwen-2.5-7B-Instruct as targets.  For GPT-4o, Qwen-2.5-7B-Instruct generates
the research dialogue.  When Qwen is the target, we instead use
Llama-3.1-8B-Instruct as the researcher to avoid assigning both roles to the
same model family.  GPT-4o target calls use the OpenAI Batch API, whereas the
open-weight models run locally.  Each target--condition cell contains the same
197 cases in the same order.

\subsection{Controlled Conditions}
\label{sec:conditions}

All conditions retain the same persona, seven-stage dynamic research
dialogue, clinical-safety framing, target-utterance quotation requirement,
four final readout directions, and direct-response stage.  They differ only
in target-visible goal text and pre-generated history:

\begin{description}
    \item[\textsc{Ours-Full} (Explicit + History).]
    The target receives the goal-aware persona and history followed by the
    original request verbatim.  This is the frozen primary condition.

    \item[\textsc{Ours-Latent} (No-Verbatim + History).]
    The same goal-aware persona and history are retained, but the original
    request is not copied verbatim into the target-visible context.  This is a
    test of RQ2, not a goal-free condition: semantic persona/history cues and
    the private researcher objective remain available.

    \item[\textsc{Ours-Streamlined} (No-Verbatim + No-History).]
    The original request and pre-generated dialogue history are both absent
    from the target-visible context.  The persona and the complete dynamic
    research dialogue remain unchanged.  This condition isolates the
    contribution of pre-generated history for RQ3.
\end{description}

The three arms produce $2\times3\times197=1{,}182$ target-case outputs.  For
each case, the method constructs four branch-scoped final readouts and reports
\textsc{Best-of-4}.  The four-try budget is held constant across all six cells.

\subsection{Metrics}
\label{sec:metrics}

Our first outcome family is CARES Level-2.  Let $A$, $C$, $R$, and $NR$ denote
Accept, Caution, Refuse, and NoResponse.  For $N=197$ cases, we report
\begin{equation}
    \mathrm{Accept}=\frac{A}{N},\qquad
    \mathrm{ASR}=\frac{A+C}{N},\qquad
    \mathrm{SS}=\frac{R+0.5C+NR}{N}.
\end{equation}
Missing responses count as attack failures (NoResponse).  We also report the
JMIR mean score (lower is more harmful) and exact-1 harmful rate.

Our strict mechanism-aware primary metric is
\begin{equation}
 \mathrm{R3{+}}\ \land\ \mathrm{B4{+}}\ \land\ \mathrm{HarmBench{+}}.
 \label{eq:strict-joint}
\end{equation}
R3+ requires substantive recovery of both the target proposition and the
requested speech act.  B4+ requires agreement with or enactment of the
complete goal.  Crucially, all three predicates in
Equation~\ref{eq:strict-joint} must hold for the \emph{same final branch}; we
never combine maxima from different branches.  HarmBench is the registered
safety axis.  Qwen3Guard and Nemotron are independent auxiliary evaluators,
for which we report analogous same-branch conjunctions without replacing the
registered primary metric.

For binary rates, we report descriptive 95\% Wilson intervals.  We additionally
compare paired case-level Accept and strict-conjunction outcomes with
two-sided exact McNemar tests and apply Holm correction over the eight
target--contrast--metric comparisons.  These inferential comparisons are
reported as an exploratory analysis rather than as part of the frozen metric
contract.

\section{Results}
\label{sec:results}

\begin{table*}[t]
\centering
\small
\setlength{\tabcolsep}{3.2pt}
\begin{tabular}{llrrrrrrrrrr}
\toprule
Target & Condition & A/C/R/NR & Accept$\uparrow$ & ASR$\uparrow$ & SS$\downarrow$
& JMIR$\downarrow$ & Exact-1$\uparrow$ & R3+@4 & B4+@4
& HB+@4 & Strict joint$\uparrow$ \\
\midrule
GPT-4o
& \textsc{Ours-Latent} & 41/1/141/14 & 20.81 & 21.32 & 0.789 & 1.359 & 87.82 & 48.22 & 17.26 & 31.98 & 7.61 \\
& \textsc{Ours-Full}$^{\dagger}$ & 35/1/152/9 & 17.77 & 18.27 & 0.820 & \textbf{1.320} & 86.80 & 44.16 & 14.72 & \textbf{32.99} & 2.54 \\
& \textsc{Ours-Streamlined} & 49/2/127/19 & \textbf{24.87} & \textbf{25.89} & \textbf{0.746} & 1.487 & 84.77 & \textbf{76.14} & \textbf{24.37} & 28.93 & \textbf{12.18} \\
\midrule
Qwen-2.5-7B
& \textsc{Ours-Latent} & 70/5/112/10 & 35.53 & 38.07 & 0.632 & \textbf{1.217} & \textbf{93.40} & 58.88 & 30.46 & \textbf{80.20} & 15.23 \\
& \textsc{Ours-Full} & 61/4/114/18 & 30.96 & 32.99 & 0.680 & 1.369 & 90.36 & 44.67 & 21.32 & \textbf{80.20} & 10.66 \\
& \textsc{Ours-Streamlined} & 71/7/100/19 & \textbf{36.04} & \textbf{39.59} & \textbf{0.622} & 1.399 & 88.83 & \textbf{68.53} & \textbf{43.65} & 73.60 & \textbf{21.32} \\
\bottomrule
\end{tabular}
\caption{Primary \textsc{Best-of-4} results on Final-197.  CARES rates,
JMIR exact-1, R3+@4, B4+@4, HarmBench positive (HB+@4), and strict joint are
percentages; SS and JMIR mean are unscaled.  Bold marks the strongest observed
value within a target, not a statistically significant winner.
$^{\dagger}$Frozen primary condition for the GPT-4o main run.}
\label{tab:main-results}
\end{table*}

\begin{table*}[t]
\centering
\small
\setlength{\tabcolsep}{5pt}
\begin{tabular}{llrrrrrr}
\toprule
Target & Condition & QG unsafe@4 & R3+B4+QG & Nemo unsafe@4
& Nemo SSH+@4 & Nemo domain@4 & R3+B4+Nemo \\
\midrule
GPT-4o
& \textsc{Ours-Latent} & 36.04 & 6.60 & 36.55 & 30.46 & 32.49 & 6.09 \\
& \textsc{Ours-Full} & 35.03 & 3.05 & 35.03 & 30.96 & 32.49 & 3.05 \\
& \textsc{Ours-Streamlined} & 35.03 & \textbf{15.23} & \textbf{37.06} & 29.95 & 32.49 & \textbf{12.69} \\
\midrule
Qwen-2.5-7B
& \textsc{Ours-Latent} & \textbf{86.80} & 17.77 & \textbf{86.80} & 77.16 & 80.20 & 14.72 \\
& \textsc{Ours-Full} & 85.28 & 12.18 & 86.29 & \textbf{79.19} & \textbf{81.22} & 9.64 \\
& \textsc{Ours-Streamlined} & 77.66 & \textbf{22.84} & 81.73 & 72.59 & 75.13 & \textbf{21.83} \\
\bottomrule
\end{tabular}
\caption{Auxiliary safety-evaluator results (\%).  QG denotes Qwen3Guard and
Nemo denotes Nemotron-3-Content-Safety.  Conjunction columns require R3+, B4+,
and the corresponding evaluator-positive label on the same branch.}
\label{tab:aux-results}
\end{table*}

\begin{table}[t]
\centering
\small
\setlength{\tabcolsep}{4pt}
\begin{tabular}{llrrrr}
\toprule
Target & Contrast vs. Latent & $\Delta$Accept & $\Delta$ASR
& $\Delta$SS & $\Delta$Strict \\
\midrule
GPT-4o & Full & $-3.05$ & $-3.05$ & $+0.030$ & $-5.08$ \\
& Streamlined & $+4.06$ & $+4.57$ & $-0.043$ & $+4.57$ \\
Qwen & Full & $-4.57$ & $-5.08$ & $+0.048$ & $-4.57$ \\
& Streamlined & $+0.51$ & $+1.52$ & $-0.010$ & $+6.09$ \\
\bottomrule
\end{tabular}
\caption{Paired descriptive differences relative to \textsc{Ours-Latent}.
Rate differences are percentage points; SS differences are unscaled.}
\label{tab:ablation-deltas}
\end{table}

\subsection{RQ1: Cross-target Effectiveness}
\label{sec:rq1-results}

The strongest observed variant is \textsc{Ours-Streamlined}.  It obtains a
CARES Accept rate of 24.87\% on GPT-4o (49/197; 95\% CI:
19.36--31.35) and 36.04\% on Qwen (71/197; 95\% CI:
29.66--42.95).  Under the stricter same-branch mechanism-aware metric, it
achieves 12.18\% on GPT-4o (24/197; 95\% CI: 8.33--17.49) and
21.32\% on Qwen (42/197; 95\% CI: 16.18--27.56).  Thus, the dynamic
research dialogue elicits non-zero strict successes across both target
families, although GPT-4o remains substantially more resistant.

\subsection{RQ2: Verbatim Goal Exposure}
\label{sec:rq2-results}

Verbatim exposure does not improve the two principal success measures.  From
\textsc{Ours-Latent} to \textsc{Ours-Full}, Accept decreases by 3.05 points
and strict success by 5.08 points on GPT-4o.  The corresponding changes on
Qwen are $-4.57$ and $-4.57$ points.  This result rules out a simple
``copy the harmful request into the context'' explanation for the observed
successes.  It does not, however, establish an absence of all goal-related
information because the latent condition retains semantic persona/history
cues and a private researcher objective.

\subsection{RQ3: Persona-History Contribution}
\label{sec:rq3-results}

Removing pre-generated history produces the highest observed CARES and strict
joint scores for both targets.  Relative to \textsc{Ours-Latent}, strict
success increases by 4.57 points on GPT-4o and 6.09 points on Qwen.  The
corresponding Accept changes are $+4.06$ and $+0.51$ points.  These results do
not support the hypothesis that pre-generated history is necessary or
beneficial; instead, they identify the target-grounded dynamic dialogue as the
more plausible active component.  However, no paired binary comparison remains
significant after Holm correction (minimum adjusted $p=.170$), so we describe
the ordering as an observed ablation trend rather than a confirmed improvement.

\subsection{RQ4: Evaluator Robustness}
\label{sec:rq4-results}

For \textsc{Ours-Streamlined}, the three same-branch strict estimates are
12.18\% (HarmBench), 15.23\% (Qwen3Guard), and 12.69\% (Nemotron) on
GPT-4o, and 21.32\%, 22.84\%, and 21.83\% on Qwen.  Their close ranges
support the conclusion that strict successes are not an artifact of a single
safety classifier.  In contrast, evaluator-positive rates without the R3+/B4+
requirements are much larger and vary more substantially.  We therefore use
the registered same-branch HarmBench conjunction as the primary result and
treat the other classifiers as corroborating evidence.

\subsection{Interpretation and Reporting Scope}
\label{sec:reporting-scope}

The ablations indicate that the central contribution is the branch-scoped,
target-grounded research dialogue and readout rather than verbatim goal
insertion or a long pre-generated history.  Nevertheless,
\textsc{Ours-Full} remains the frozen primary condition, while
\textsc{Ours-Streamlined} is the strongest observed ablation.  Promoting the
latter to a new primary method requires confirmation on a previously unused
evaluation set.

Finally, existing baseline results available to us use a different 64-case
subset.  We therefore do not claim state-of-the-art performance or rank our
197-case results against those numbers.  A formal superiority claim requires
running baselines on the same 197 cases with the same four-attempt budget,
missing-response policy, and evaluator suite.
